\section{Discussion}
\label{sec:discussion}

The effect of local epochs on federated Lasso selection depends on the
aggregation procedure and penalty rule. The orthogonal calculation makes
one distinction exact: local CD finishes in one sweep, so further local
epochs have no effect even though averaging can enlarge the selected support.
For general designs, an averaged fixed point need not minimise the pooled
Lasso objective. Neither fact implies that every average is dense or that
its statistical performance is determined by its optimisation gap alone.

The simulation comparisons support reporting tuning choices alongside
aggregation choices. Local one-SE penalties improve plain averaging, while
post-fit and per-round thresholds introduce additional selection steps.
Their search spaces differ, and the resulting performance differences
should not be assigned wholly to one mechanism. Likewise, favourable
support recovery at an inherited penalty should not be interpreted as a
universal advantage of a particular solver.

Activity thresholds, numbers selected, and variable identities are distinct
reporting choices. Sensitivity to the numerical activity threshold should
be shown, and a nearly constant number of selected variables should not be
called stable support without checking the identities. The fixed synthetic
benchmark illustrates this distinction: thresholding makes the size curve
flatter while its agreement with a pooled reference varies more with $E$.
Where a generating support is unavailable, such agreement remains a
descriptive comparison between estimators.

Communication accounting must include additional federated fits required
for tuning. Post-fit thresholding is inexpensive under the stated upload
convention; selecting a per-round threshold requires several complete
trajectories. Local CD work and uploaded bytes can move in opposite
directions as $E$ changes. Neither measure alone establishes a preferred
schedule, especially when different algorithm families perform different
local computations or use different stopping criteria.

\paragraph{Limitations.}
The exact support calculations require orthogonal local designs with
$p\le\min_jm_j$, outside the high-dimensional simulations. The
general-design fixed-point result supplies no convergence rate or
nonasymptotic bound linking $E$ to support error. The main study uses one
sample size, one sparsity level, three sites, Gaussian designs and a shared
coefficient vector. It does not establish uniform performance over other
site counts, signal strengths, coefficient heterogeneity or response models.
The observation-level one-SE rule is a tuning heuristic, especially with
heterogeneous fixed sites; it is not uncertainty quantification for the
selected support. A systematic comparison with stratified or nested tuning
is outside this study.

The original fixed benchmark cannot validate recovery because its generating
coefficients and algorithm are unavailable. The public-data illustration
uses simulated site partitions and does not represent a deployed
multi-institutional analysis. Repeated splits describe sensitivity within
that finite dataset, not independent population replications. The public example also includes capped, unconverged averaging iterates,
which limits their interpretation as fixed points. The
communication-budget comparison is a deterministic specialization of
FedDualAvg; it does not reproduce every stochastic setting or accelerated
method in the composite federated literature. Local computation and tuning
burden remain relevant to that comparison.

Finally, keeping records at a site does not provide a formal privacy
guarantee. Parameters and scalar summaries may reveal sample information
\citep{zhu2021}. The experiments measure algorithmic and statistical
behaviour, not differential privacy, secure aggregation, or operational
network performance.
